% Appendix for the Nera draft (Oct 22-23, 2026). The additional-policy table (7.9 point) and the CES appendix
% (14.40 calibration) of the mock are left out: both predate the current calibration and supply rule.
\section{Household Optimality Conditions}
\label{app:household}

In this Appendix, I derive the conditions that Section~\ref{sec:m-analysis}
uses: the choices within a housing option, the Euler equation for the bond
and the borrowing multiplier $\zeta_t$, the effect of credit on the value of
a household and on fertility, and the two results in the footnotes of that
Section. Throughout, $n$ is the number of children at home after the birth
outcome of period $t$, and the budget constraint \eqref{eq:m-budget} of the
chosen option defines $c_t$ as a function of $b_{t+1}$ and, for a renter, of
$h^R_t$.

\subsubsection*{Choices within an option.}
Within option $j$ of \eqref{eq:m-vh}, renting or owning $h^O_t\in\mathcal H$,
the household chooses $b_{t+1}$ and, if it rents, $h^R_t$, subject to the
borrowing limit \eqref{eq:m-ltv}, with $h^O_t=0$ for a renter, and to the
cap $h^R_t\le\bar h^R$. Let $\zeta_t\ge0$ and $\nu_t\ge0$ be the multipliers
on the two constraints. By \eqref{eq:m-budget},
$\partial c_t/\partial b_{t+1}=-1$ in every option and
$\partial c_t/\partial h^R_t=-r_t$ for a renter. In \eqref{eq:m-utility}, the
marginal utilities of consumption and of rented space are the same power of
the composite in brackets, times $\alpha_0/c_t$ and
$(1-\alpha_0)/(h^R_t-\bar h(n))$, so the condition
$\partial u/\partial h^R_t=r_t\,\partial u/\partial c_t+\nu_t$ gives, below
the cap,
\begin{equation}
 r_t\bigl(h^R_t-\bar h(n)\bigr)=\frac{1-\alpha_0}{\alpha_0}\,c_t .
 \label{eq:app-intra}
\end{equation}
A renter spends the share $\alpha_0$ of $c_t+r_t(h^R_t-\bar h(n))$ on
consumption and $1-\alpha_0$ on space above the requirement. With
$b_t=b_{t+1}=0$ and no home to sell, this spending is
$y_t+T_t-r_t\bar h(n)$, and substituting the two shares into
\eqref{eq:m-utility} gives \eqref{eq:m-indirect}. At the cap,
$h^R_t=\bar h^R$ and
$\nu_t=\partial u/\partial h^R_t-r_t\,\partial u/\partial c_t\ge0$. An
owner's space is fixed within the option, so its only choice is the bond.

\subsubsection*{The bond.}
The first-order condition for $b_{t+1}$ is the Euler equation
\begin{equation}
 \frac{\partial u}{\partial c_t}
 =\beta\,\E_t\Bigl[s_a\frac{\partial V_{t+1}}{\partial b_{t+1}}+(1-s_a)\,B'(x_{t+1})\Bigr]+\zeta_t,
 \qquad
 \zeta_t\bigl(b_{t+1}+\phi P_th^O_t\bigr)=0,
 \label{eq:app-euler}
\end{equation}
where $h^O_t=0$ for a renter and
$B'(x_{t+1})=\theta_0(\theta_1+x_{t+1})^{-\sigma}$. The multiplier $\zeta_t$
is the gain, in period-$t$ utility, from borrowing one more unit in period
$t$ and repaying $R_b$ in period $t+1$: the marginal utility of consumption
today net of the discounted expected marginal value of the bond tomorrow. It
is positive only at the limit. The estate $x_{t+1}=b_{t+1}+(1-\mu)P_th^O_t$
contains the bond position itself, so $\partial x_{t+1}/\partial b_{t+1}=1$
and the bequest term carries no $R_b$; the survival term does, through the
budget of period $t+1$.

For Gumbel shocks with scale $\kappa$, the expected maximum over options is
the logsum in \eqref{eq:m-fertility} and \eqref{eq:m-tenure}, whose
derivative with respect to the value of one option is the probability of
that option:
\begin{equation}
 \frac{\partial}{\partial V^k}\,\kappa\log\sum_{j}\exp\bigl(V^j/\kappa\bigr)
 =\frac{\exp\bigl(V^k/\kappa\bigr)}{\sum_{j}\exp\bigl(V^j/\kappa\bigr)} .
 \label{eq:app-logsum}
\end{equation}
The household observes the shocks before it chooses, so a marginal change in
a state moves the value of the chosen option only; at the shocks where the
choice switches, the two options have the same value. The derivative of the
expected maximum with respect to a state is therefore the
probability-weighted average of the derivatives of the option values. When
$\sigma\ge1$, the value of an option whose feasible set shrinks to a point
tends to $-\infty$, since $u\to-\infty$ as $c_t\to0$, so its probability
tends to zero and the logsum is continuous where an option enters or leaves
the feasible set. Within an option $j$ of period $t+1$, $b_{t+1}$ enters
only the budget, with $\partial c_{t+1}/\partial b_{t+1}=R_b$, so
$\partial V^j_{t+1}/\partial b_{t+1}=R_b\,\partial u/\partial c_{t+1}$ at the
option's consumption, whether or not a constraint binds in that period.
Writing $\frac{\partial u}{\partial c_{t+1}}(j,n')$ for this marginal
utility with $n'$ children at home, and applying \eqref{eq:app-logsum} to
\eqref{eq:m-tenure} and to \eqref{eq:m-fertility},
\begin{multline}
 \frac{\partial V_{t+1}}{\partial b_{t+1}}
 =R_b\Bigl[\bigl(1-\pi_{a+1}\pi^F_{t+1,a+1}\bigr)\sum_{j}\pi^H_{t+1}(j\mid n_{t+1})\,\frac{\partial u}{\partial c_{t+1}}(j,n_{t+1})\\
 +\pi_{a+1}\pi^F_{t+1,a+1}\sum_{j}\pi^H_{t+1}(j\mid n_{t+1}+1)\,\frac{\partial u}{\partial c_{t+1}}(j,n_{t+1}+1)\Bigr],
 \label{eq:app-envelope}
\end{multline}
where $n_{t+1}$ is the number of children at home at the start of period
$t+1$ and $\pi_{a+1}\pi^F_{t+1,a+1}$ is the probability of a birth in that
period, zero outside the fertile ages. The derivative is $R_b$ times the
marginal utility of consumption in period $t+1$, averaged over the birth
outcome and the housing choice of that period; the cost $\xi$ drops out. The
expectation $\E_t$ in \eqref{eq:app-euler} then averages over $z_{t+1}$ and
the departure of children.

\subsubsection*{A birth and the multiplier.}
At the limit, $\zeta_t$ is the difference between the two sides of
\eqref{eq:app-euler} at $b_{t+1}=-\phi P_th^O_t$, and a birth changes both.
In period $t$, it adds $r_th_P$ to the spending a renter cannot avoid,
lowers an owner's space above the requirement by $\chi h_P$, and raises
$e(n)$; for $\sigma>1$, each of these raises $\partial u/\partial c_t$ at the
given bond position. In period $t+1$ the same holds with probability
$1-\lambda$, when the child is still at home, and the bequest term does not
depend on $n$. Within an option,
\begin{equation}
 \zeta_t(h^O_t,n+1)-\zeta_t(h^O_t,n)
 =\Bigl[\frac{\partial u}{\partial c_t}\Bigr]_{n}^{n+1}
 -\beta s_a\,\E_t\Bigl[\frac{\partial V_{t+1}}{\partial b_{t+1}}\Bigr]_{n}^{n+1},
 \label{eq:app-birth}
\end{equation}
where $[\,\cdot\,]_n^{n+1}$ is the difference between the values with $n+1$
and with $n$ children at home, at $b_{t+1}=-\phi P_th^O_t$ in both. The
first term is positive, and larger when spending in period $t$ is small,
since the marginal utility of consumption is convex in spending. The second
is scaled by $\beta s_a$ and by $1-\lambda$, averaged with the choice
probabilities of period $t+1$, and evaluated at that period's spending,
which is larger when earnings rise with age. The multiplier therefore rises
for a young household with low spending today and rising earnings, and can
fall for one that expects to spend less next period; the sign is ambiguous
in general.

\subsubsection*{Credit.}
Consider a change in $\phi$ in period $t$ alone. It enters the problem of an
owner only through the borrowing limit, whose derivative with respect to
$\phi$ is $P_th^O_t$, so by the envelope theorem
$\partial V^j_t/\partial\phi=\zeta_t(h^O,n)P_th^O$ for owning
$h^O\in\mathcal H$ and $\partial V^j_t/\partial\phi=0$ for renting, where
$\zeta_t(h^O,n)$ is the multiplier in the option of owning $h^O$ with $n$
children at home. Applying \eqref{eq:app-logsum} to \eqref{eq:m-tenure}
gives the second line of \eqref{eq:m-credit}. For the probability of trying,
\eqref{eq:m-fertility} and
$V^F_{t,A}-V^F_{t,0}=\pi_a\bigl[V^H_t(n_t+1)-V^H_t(n_t)-\xi\,\1\{\text{first child}\}\bigr]$
give
\begin{equation}
 \frac{\partial\pi^F_{t,a}}{\partial\phi}
 =\frac{\pi^F_{t,a}\bigl(1-\pi^F_{t,a}\bigr)}{\kappa}\,
 \frac{\partial\bigl(V^F_{t,A}-V^F_{t,0}\bigr)}{\partial\phi}
 =\frac{\pi^F_{t,a}\bigl(1-\pi^F_{t,a}\bigr)}{\kappa}\,\pi_a
 \left[\frac{\partial V^H_t(n_t+1)}{\partial\phi}-\frac{\partial V^H_t(n_t)}{\partial\phi}\right],
 \label{eq:app-credit-pi}
\end{equation}
where $\kappa=\kappa_1$ before the first child and $\kappa_C$ after; the
cost $\xi$ does not depend on $\phi$ and drops out. This is the first line
of \eqref{eq:m-credit}. If the change in $\phi$ is permanent, it also enters
$V_{t+1}$, and the envelope theorem adds
$\beta s_a\E_t\,\partial V_{t+1}/\partial\phi$ to the derivative of each
option, which by \eqref{eq:app-logsum} is the second line of
\eqref{eq:m-credit} in period $t+1$, averaged over the choices of that
period, and so on: the derivative becomes the sum of these terms over all
future periods, discounted by $\beta$ and weighted by survival and by the
probabilities of the choices that lead to each option.

The bracket in \eqref{eq:m-credit} splits into two parts,
\begin{multline}
 \frac{\partial V^H_t(n+1)}{\partial\phi}-\frac{\partial V^H_t(n)}{\partial\phi}
 =\sum_{h^O\in\mathcal H}\bigl[\pi^H_t(h^O\mid n+1)-\pi^H_t(h^O\mid n)\bigr]\,\zeta_t(h^O,n)\,P_th^O\\
 +\sum_{h^O\in\mathcal H}\pi^H_t(h^O\mid n+1)\bigl[\zeta_t(h^O,n+1)-\zeta_t(h^O,n)\bigr]\,P_th^O ,
 \label{eq:app-credit-decomp}
\end{multline}
where the first sum holds the multipliers at their values without the child
and the second holds the probabilities at their values with the child. The
first sum is positive when the child moves probability from renting, and
from homes where the limit is slack, toward homes where
$\zeta_t(h^O,n)P_th^O$ is large. The second is positive when the child
raises the multiplier in the homes the household is likely to own with it,
with the sign of \eqref{eq:app-birth}. Both are small for a household that
rents with high probability with and without the child.

\subsubsection*{The owner's cost of space.}
By \eqref{eq:m-budget}, an owner that holds a unit of housing from period $t$
to period $t+1$ pays $P_t$ and $\tau^pP_t$ in period $t$. In period $t+1$ it
pays maintenance $\delta_HP_{t+1}$ if it stays, or receives
$(1-\delta_H-\mu)P_{t+1}$ if it sells, so apart from the selling cost the unit
is worth $(1-\delta_H)P_{t+1}$ in period $t+1$, which the bond values at
$(1-\delta_H)P_{t+1}/R_b$ in period $t$. A renter pays $r_t$ in period $t$. By
\eqref{eq:m-rent},
\begin{equation}
 r_t-\Bigl[(1+\tau^p)P_t-\frac{(1-\delta_H)P_{t+1}}{R_b}\Bigr]=0,
 \label{eq:app-usercost}
\end{equation}
where the term in brackets is the owner's cost in period-$t$ units. Apart from
the selling cost, owning a unit of space costs the same as renting it, along a
transition as well as in a stationary equilibrium.

\subsubsection*{The cost of a child's space and earnings.}
Let $\tilde y_t=y_t+T_t$, $k=r_th_P$ and
$\Theta_t=\alpha_0^{\alpha_0}(1-\alpha_0)^{1-\alpha_0}r_t^{\alpha_0-1}$. The
first term of \eqref{eq:m-indirect} is $(\Theta_t\tilde y_t)^{1-\sigma}/(1-\sigma)$
without children and
$\bigl(\Theta_t(\tilde y_t-k)/e(1)\bigr)^{1-\sigma}/(1-\sigma)$ with one
child, so the utility cost of a first child and its second derivative are
\begin{equation}
\begin{aligned}
 \Delta_t(\tilde y_t)&=\frac{\Theta_t^{1-\sigma}}{1-\sigma}\Bigl[\tilde y_t^{\,1-\sigma}-e(1)^{\sigma-1}(\tilde y_t-k)^{1-\sigma}\Bigr],\\
 \Delta_t''(\tilde y_t)&=\sigma\,\Theta_t^{1-\sigma}\Bigl[e(1)^{\sigma-1}(\tilde y_t-k)^{-\sigma-1}-\tilde y_t^{\,-\sigma-1}\Bigr],
\end{aligned}
\label{eq:app-delta}
\end{equation}
where $\tilde y_t>k$. For $\sigma\ge1$, $e(1)^{\sigma-1}\ge1$ and
$(\tilde y_t-k)^{-\sigma}>\tilde y_t^{\,-\sigma}$, so
$\Delta_t'(\tilde y_t)=\Theta_t^{1-\sigma}\bigl[\tilde y_t^{\,-\sigma}-e(1)^{\sigma-1}(\tilde y_t-k)^{-\sigma}\bigr]<0$
and $\Delta_t''>0$: the cost is decreasing and convex in earnings. The
second derivative rises with $k$, since
$\partial\Delta_t''/\partial k=\sigma(\sigma+1)\Theta_t^{1-\sigma}e(1)^{\sigma-1}(\tilde y_t-k)^{-\sigma-2}>0$,
and, for $\sigma>1$, with $r_t$, since $\Theta_t^{1-\sigma}$ is proportional
to $r_t^{(1-\alpha_0)(\sigma-1)}$. For a renter at the cap with and without
the child, $c_t=\tilde y_t-r_t\bar h^R$ in both cases and the child lowers
space above the requirement from $\bar h^R$ to $\bar h^R-h_P$, so
\[
 \Delta_t(\tilde y_t)=\frac{(\tilde y_t-r_t\bar h^R)^{\alpha_0(1-\sigma)}}{1-\sigma}
 \Bigl[(\bar h^R)^{(1-\alpha_0)(1-\sigma)}-e(1)^{\sigma-1}\bigl(\bar h^R-h_P\bigr)^{(1-\alpha_0)(1-\sigma)}\Bigr],
\]
which for $\sigma>1$ is positive, decreasing and convex in $\tilde y_t$ and
rises with $r_t$ and $h_P$.

\subsubsection*{The figures of Section~\ref{sec:m-analysis}.}
\label{app:model-figures}
I compute Figures~\ref{fig:m-risk} and~\ref{fig:m-illustration} from the model at the parameters of
Table~\ref{tab:quant-parameters}. The estimated values are
$\beta_{\text{ann}}=0.9357$, $\psi_0=0.1774$, $\gamma=0.6457$, $\xi=0.2946$,
$\kappa_1=0.0505$, $\kappa_C=0.1285$, $h_P=0.8899$, $\chi=1.1121$,
$\kappa_T=0.0131$ and $\theta_0=0.2435$, and the assigned parameters are those
of the same table. In the 2007 stationary equilibrium at these values, the
price is $P_0=0.7707$, the rebate is $T=0.1879$ per household and there are
$N_0=1$ households. Supply \eqref{eq:m-supply} is $H=6.776\,P^{0.63}$. Each
solve below is a stationary solve of the household problem and of the
distribution of households at a given price, with the rebate set so that the
government budget balances at that price (relative residual below
$10^{-6}$). None of these solves is a transition. The stationary rent is
$r=\bigl[1+\tau^p-(1-\delta_H)/R_b\bigr]P=0.1698\,P$. The axes of the figures
carry no numerical scale, and I report the values of the marked points
below.

\paragraph{Figure~\ref{fig:m-risk}.} This figure is computed from
\eqref{eq:m-indirect} and involves no solve. With $\sigma=2$, the cost of a
first child is
$\Delta(y)=\bigl[e(1)/(y+T-rh_P)-1/(y+T)\bigr]/\Theta$, where
$\Theta=\alpha_0^{\alpha_0}(1-\alpha_0)^{1-\alpha_0}r^{\alpha_0-1}$. The
values are $\alpha_0=0.733$, $e(1)=1.35^{0.7}=1.234$ and $h_P=0.890$, and
both curves use the 2007 rebate, $T=0.188$. The solid curve uses the 2007
rent, $r=0.1308$, so that $rh_P=0.116$. The dashed curve uses the rent at
the long-run price $P_1=0.394$ of Figure~\ref{fig:m-illustration},
$r=0.0669$, so that $rh_P=0.060$. Earnings $y_L=0.274$ and $y_H=0.752$ are
after-tax earnings per period at ages 22--25 in the first and third of the
nine efficiency states, $z=0.103$ and $z=0.284$. In the 2007 solve, at least
98\% of the renters of that age in these two states end the period at the
borrowing limit. Over the plotted range, neither a childless renter nor a
parent reaches the rental cap $\bar h^R=6$ at either rent. The additional
cost due to risk, $\tfrac12[\Delta(y_L)+\Delta(y_H)]-\Delta(\bar y)$, is 0.245
at the 2007 rent and 0.106 at the long-run rent; the braces mark these two
differences.

\paragraph{Figure~\ref{fig:m-illustration}.} Both panels come from 22 solves at
given prices. Ten are at $\psi_0=0.1774$, at prices from 0.30 to 0.92. The
other twelve are at $\psi_1=0.1472$ over the same range, three of them steps
of the search for $P_1$. Here $\psi_1$ is the long-run level $\psi'$ of the
decay path \eqref{eq:psi-decay}, with $\rho=0.35$.
\textcolor{red}{[$\psi_1$ is from a nearby point; refit at the final
calibration.]} At each price, $n(P;\psi)$ is births per entering household,
with the bin of three or more children counted as 3.60 and replacement at
$\bar n=2.1$, and $h(P;\psi)$ is housing demand per household. The demand and
fertility curves interpolate the solved points with monotone cubic splines,
and supply is $H=6.776\,P^{0.63}$. At $A$, $P_0=0.771$, $H=5.750$, $N_0=1$
and $n=\bar n$. At $A'$, preferences are $\psi_1$ at $P_0$ with $N_0$
households, and $n=1.632$. At $B$, the long-run stationary equilibrium,
$P_1=0.394$, $n=\bar n$ (births within $10^{-6}$ of replacement), $H=3.769$,
$N_1=H/h(P_1;\psi_1)=0.472$ and $T=0.133$. From $A$ to $B$, the price falls by
49\%, the stock by 34\% and the number of households by 53\%. In panel (a),
demand at $A'$ is $N_0h(P_0;\psi_1)=5.698$, which is 0.9\% below demand at
$A$, so I mark $A'$ at $A$. The arrows only indicate the direction of the
adjustment.

% Slow-stock (KMV permit) policy results, moved from Section 6 on Oct 22, 2026 (his decision Oct 21 22:36Z: free entry
% in the main text, KMV in the appendix as robustness). Numbers: output/model/experiments/housing_supply/20261021_paper_results/
% p6_values.json (p6_figs.py), which match 20261017_accounting/report_20261018/policy_runs_table_f3.md (f3 runs, corrected
% 2023 start). Matched supply at eta 0.63 never converged, so it is left out here.
\providecommand{\pend}{\textcolor{red}{[pending]}}
\providecommand{\PRF}{../../output/model/experiments/housing_supply/20261021_paper_results}

\section{Robustness of the Policy Results}
\label{app:policy-robust}

This appendix repeats the experiments of Section~\ref{sec:policy} with a
slowly adjusting stock, at two supply elasticities.

\subsection{A slowly adjusting stock}
\label{app:policy-slow}

Here the housing stock adjusts slowly, as in
\citet{kaplanMitmanViolanteHousing2020}. Builders combine goods $X_t$ with a
fixed flow of permits $\bar L$ each period, $I_t=X_t^{\alpha}\bar L^{1-\alpha}$,
so the marginal cost of a unit rises with this period's construction, not with
the stock. The permit rent, $(1-\alpha)P_tI_t$, goes to the government, which
spends it on goods households do not value. Setting $\alpha/(1-\alpha)=\eta$
and the scale of construction to $\delta_HH_0$, construction and the stock
follow
\begin{equation}
 I_t=\delta_H H_0P_t^{\eta},\qquad H_t=(1-\delta_H)H_{t-1}+I_t .
 \label{eq:app-slow-stock}
\end{equation}
In a stationary equilibrium construction replaces depreciation, so the stock
lies on the same curve $H=H_0P^{\eta}$ as under free entry, and the 2007
equilibrium is the same. Along a transition only construction responds to the
price within a period, so the stock responds with an elasticity of about
$\delta_H\eta$, and the price absorbs most of a change in demand. The decline
in the preference for children from 2007 to 2023 is fitted under this rule, so
the 2023 state differs from the one in Section~\ref{sec:policy}. Everything
else is as in Section~\ref{sec:policy-design}.

Table~\ref{tab:app-policy-slow} reports the results. With $\eta=0.63$,
mortgage credit raises births by 1.7\% in 2023, but the gain falls to 0.6\%
by 2063, and the completed fertility of the 2023 cohort falls by 0.2\%.
Ownership at ages 18 to 29 rises by 19 percentage points on impact. A 10\%
supply expansion lowers the price by 3.0\% in 2023 and raises births by only
0.4\%. Construction keeps adding space, rents are 4\% lower by 2063, births
are 3.2\% higher, and the completed fertility of the 2023 cohort rises by
1.4\%. Credit for parents raises births by 6.2\% in 2023 and by 17\% by 2063,
and the completed fertility of the 2023 cohort by 5.9\%. As the additional
families demand more space, the house price rises, by 9\% by 2063.
Figure~\ref{fig:app-policy-slow-paths} shows the paths.

With $\eta=1.5$ (panel B), the effects on births in 2023 and 2063 are close
to those with $\eta=0.63$. To raise births in 2023 as much as mortgage credit,
the supply expansion must be 62\%. It then raises births by 12\% by 2063 and
the completed fertility of the 2023 cohort by 5.6\%, against $-0.8\%$ under
credit.

\problem{With $\eta=1.5$ the paths are still well short of the terminal
steady state in 2403.}

\begin{table}[H]
\centering
\caption{Housing policies from 2023 with a slowly adjusting stock}
\label{tab:app-policy-slow}
\small
\begin{tabularx}{\textwidth}{@{}Xrrrcrr@{}}
\toprule
 & \multicolumn{3}{c}{Births} & Completed fertility & \multicolumn{2}{c}{House price} \\
\cmidrule(lr){2-4}\cmidrule(l){6-7}
 & 2023 & 2063 & 2103 & 2023 cohort & 2023 & 2063 \\
\midrule
\multicolumn{7}{@{}l}{\emph{A. Supply elasticity $\eta=0.63$}}\\
Mortgage credit & $+1.7$ & $+0.6$ & $+0.9$ & $-0.2$ & $+0.2$ & $+0.2$ \\
Housing supply, $+10\%$ & $+0.4$ & $+3.2$ & $+5.1$ & $+1.4$ & $-3.0$ & $-4.7$ \\
Credit for parents & $+6.2$ & $+17.1$ & $+23.1$ & $+5.9$ & $+1.8$ & $+9.4$ \\
\addlinespace
\multicolumn{7}{@{}l}{\emph{B. Supply elasticity $\eta=1.5$}}\\
Mortgage credit & $+1.9$ & $+0.6$ & $+0.0$ & $-0.8$ & $-0.1$ & $+0.2$ \\
Housing supply, $+10\%$ & $+0.3$ & $+2.4$ & $+3.8$ & $+1.1$ & $-2.1$ & $-3.1$ \\
Matched supply, $+62\%$ & $+1.9$ & $+12.4$ & $+20.6$ & $+5.6$ & $-11.1$ & $-15.2$ \\
Credit for parents & $+6.4$ & $+17.2$ & $+25.6$ & $+7.5$ & $+1.1$ & $+6.0$ \\
\bottomrule
\end{tabularx}
\par\smallskip
\begin{minipage}{0.95\linewidth}
\footnotesize\textit{Notes:} \% differences from the no-policy path from the
same 2023 state, at the same supply elasticity, with the stock following
\eqref{eq:app-slow-stock}. Births: births in the four-year period starting in
the year shown. Completed fertility: children ever born to the cohort aged
18--21 in 2023. Without policy and with $\eta=0.63$, period fertility is 1.64
in 2023, 1.80 in 2063 and 2.00 in 2103. Matched supply: the supply shift that
raises births in 2023 as much as mortgage credit.
\end{minipage}
\end{table}

\begin{figure}[H]
\centering
\includegraphics[width=0.9\linewidth]{\PRF/fig_p6_paths_eta063.pdf}
\caption{Mortgage credit, housing supply and credit for parents with a slowly
adjusting stock}
\label{fig:app-policy-slow-paths}
\par\smallskip
\begin{minipage}{0.92\linewidth}
\footnotesize\textit{Notes:} \% differences from the no-policy path from the
same 2023 state, $\eta=0.63$, stock following \eqref{eq:app-slow-stock}.
Births in each four-year period (panel a), completed fertility of the cohort
aged 18--21 in the year shown (panel b), number of households (panel c) and
house price (panel d). Housing supply: the supply curve shifted out by 10\%.
\end{minipage}
\end{figure}


\section{Solution Algorithms}
\label{app:algorithms}

\paragraph{Stationary equilibrium.} I solve for the house price $P$ and the
rebate $T$, given parameters, a fixed preference for children and the payroll
tax. The pension is set before the search: in a stationary population the
age distribution follows from survival and entry, so the pension that
balances the payroll tax is known in advance, and the balance is verified at
the end.
\begin{enumerate}
\item Guess $P$ and derive the rent $r$ from the user cost.
\item Solve the household problem backward over age and compute the choice
probabilities.
\item Construct the stationary distribution of households, iterating on the
rebate until $T$ times the number of households equals property-tax revenue.
In the calibration the number of households is one and $H_0$ is set so that
supply equals demand at $P$; in the counterfactuals $H_0$ is fixed and the
number of households equals housing supply divided by mean housing per
household.
\item Compute the renewal residual: adjusted births divided by 2.1, relative to the
number of entrants, minus one. It is zero when the population reproduces
itself.
\item Update $P$ by bracketing the root of the renewal residual and repeat
until the renewal and rebate residuals are below $10^{-6}$. Re-solve at the
selected price and require the same residuals, and verify that the pension
budget balances to the same tolerance.
\end{enumerate}
% [CODE] native_phase_b.py 25-26 (RENEWAL_TOL, PAYGO_TOL 1e-6), 97-99
% (supply rule), 126-127 (population = supply/demand; renewal = births/(2.1
% entry) - 1), 440-454 (bracket, selected root, exact repeat);
% equilibrium.py (model) 84-96 (implied H0 under population_one, renewal
% root); fiscal_closure.py (T N = revenue - grants, FISCAL_TOL 1e-6);
% e5f_stationary_paygo.py 47-62 (pension from the payroll tax). The earlier
% text's "excess demand for housing" residual and one-period reproduction
% check are not in the engine.

\paragraph{Transition.} I solve for the paths $(P_t,m_t,T_t)_{t\geq t_0}$,
given the inherited households and a preference believed permanent.
\begin{enumerate}
\item Solve the new stationary equilibrium; it supplies the terminal value
function.
\item Guess the three paths and derive rents from \eqref{eq:m-rent}.
\item Solve the household problem backward in calendar time from the terminal
value.
\item Advance households and the queue of entrants forward from the inherited
state, and aggregate housing demand and taxes at each date.
\item Compute the three residuals at every date. The loss is the largest
residual across dates.
\item Update the paths jointly and repeat until the loss is below tolerance.
Extend the horizon to check that early responses are stable.
\end{enumerate}
The transitions in the paper use 100 four-year periods, from 2007 (or 2023
for the policy experiments) to 2403, with tolerances of $2\times10^{-4}$ on
the housing market and $2\times10^{-5}$ on the pension and the rebate, and
$10^{-6}$ on the renewal condition of the terminal steady state.
% [CODE] gates from state_2023/pair_result.json (market 2e-4, fiscal 2e-5,
% stationary renewal 1e-6, horizon 1e-3); 7p9167_v6/table.md notes for the
% LR_PHI100 and LR_F_PHI100 endpoints (renewal gap <= 0.003).
